% Preamble для ShopApp-Book — подключается через --include-in-header в build-pdf.sh
% Этот файл — чистый LaTeX. pandoc передаёт его в xelatex без изменений.
% YAML-блок header-includes: | в metadata.yaml ломал комментарии (% → \%)
% и склеивал строки, поэтому критичные декларации перенесены сюда.

\usepackage{fancyhdr}
\pagestyle{fancy}
\fancyhf{}
\fancyhead[L]{\small\nouppercase\leftmark}
\fancyhead[R]{\small\thepage}
\renewcommand{\headrulewidth}{0.3pt}

\usepackage{xcolor}
\definecolor{quotecol}{HTML}{4A4A4A}
\usepackage{microtype}
\raggedbottom
\clubpenalty=10000
\widowpenalty=10000
% Без \thechapter — номер главы уже в тексте заголовка («Глава N. ...»)
\renewcommand\chaptermark[1]{\markboth{#1}{}}

\usepackage{tocloft}
\setlength{\cftchapnumwidth}{2.4em}
\setlength{\cftsecnumwidth}{3.2em}
\setlength{\cftsubsecnumwidth}{4em}
\setlength{\cftsubsubsecnumwidth}{4.8em}

\let\oldchapter\chapter
\renewcommand{\chapter}{\clearpage\oldchapter}

\usepackage{enumitem}
\setlistdepth{9}
\renewlist{itemize}{itemize}{9}
\setlist[itemize,1]{label=$\bullet$}
\setlist[itemize,2]{label=$\circ$}
\setlist[itemize,3]{label=$\ast$}
\setlist[itemize,4]{label=$\diamond$}
\setlist[itemize,5]{label=$\cdot$}
\setlist[itemize,6]{label=$\bullet$}
\setlist[itemize,7]{label=$\circ$}
\setlist[itemize,8]{label=$\ast$}
\setlist[itemize,9]{label=$\diamond$}
\renewlist{enumerate}{enumerate}{9}
\setlist[enumerate,1]{label=\arabic*.}
\setlist[enumerate,2]{label=\alph*.}
\setlist[enumerate,3]{label=\roman*.}
\setlist[enumerate,4]{label=\arabic*.}
\setlist[enumerate,5]{label=\alph*.}

\usepackage{newunicodechar}
\usepackage{amssymb}
\usepackage{textcomp}

% Apple Symbols + Menlo для хоткеев macOS
\newfontfamily\applesymbols{Apple Symbols}
\newfontfamily\menlosym{Menlo}

% Математические стрелки и операторы — \ensuremath работает в любом контексте
% (внутри \textbf, \texttt, \href и т.д.).
\newunicodechar{→}{\ensuremath{\rightarrow}}
\newunicodechar{←}{\ensuremath{\leftarrow}}
\newunicodechar{↑}{\ensuremath{\uparrow}}
\newunicodechar{↓}{\ensuremath{\downarrow}}
\newunicodechar{⇒}{\ensuremath{\Rightarrow}}
\newunicodechar{⇐}{\ensuremath{\Leftarrow}}
\newunicodechar{≈}{\ensuremath{\approx}}
\newunicodechar{≠}{\ensuremath{\neq}}
\newunicodechar{≤}{\ensuremath{\leq}}
\newunicodechar{≥}{\ensuremath{\geq}}
\newunicodechar{·}{\ensuremath{\cdot}}
\newunicodechar{…}{\ldots}

% Чек-маркеры
\newunicodechar{✓}{\ensuremath{\checkmark}}
\newunicodechar{✗}{\ensuremath{\times}}
\newunicodechar{●}{\ensuremath{\bullet}}
\newunicodechar{○}{\ensuremath{\circ}}

% Хоткеи macOS
\newunicodechar{⌘}{{\applesymbols ⌘}}
\newunicodechar{⌥}{{\applesymbols ⌥}}
\newunicodechar{⌃}{{\applesymbols ⌃}}
\newunicodechar{↩}{{\applesymbols ↩}}
\newunicodechar{⎋}{{\applesymbols ⎋}}
\newunicodechar{⌫}{{\applesymbols ⌫}}
\newunicodechar{⏎}{{\applesymbols ⏎}}

% ⇧ U+21E7 не входит в Apple Symbols — берём из Menlo
\newunicodechar{⇧}{{\menlosym ⇧}}
\newunicodechar{▶}{{\menlosym ▶}}
\newunicodechar{■}{{\menlosym ■}}

% Эмодзи в заголовках книги. У xelatex/tectonic нет встроенного emoji-фонта,
% поэтому заменяем на текстовые маркеры.
\newunicodechar{📋}{[ITOG]\,}
\newunicodechar{🎉}{[!]\,}
\newunicodechar{🛠}{[TASK]\,}
\newunicodechar{💡}{[IDEA]\,}
\newunicodechar{⚠}{[!]\,}
\newunicodechar{🚧}{[WIP]\,}
\newunicodechar{☕}{[coffee]\,}
\newunicodechar{⏳}{[..]\,}
\newunicodechar{✅}{\ensuremath{\checkmark}\,}
% Валютные символы — нет в Menlo Italic.
\newunicodechar{₸}{T}
\newunicodechar{₽}{R}
\newunicodechar{₾}{L}
\newunicodechar{֏}{D}
